\documentclass[10pt,letterpaper]{article}
\usepackage{cornell-notes}

\title{Clause 33.06.07.02: Function Call Once}
\author{}
\date{}
\setDocTitle{Clause 33.06.07.02: Function Call Once}
\setDocSubtitle{C++ 2024}
\setDocOwner{C++ 2024 Cornell Notes}
\setDocDate{\today}
\setCornellCollection{C++ 2024 Cornell Notes}
\setCornellUnitType{Clause}
\setCornellUnitNumber{33.06.07.02}
\setCornellUnitTitle{Function Call Once}
\hypersetup{pdftitle={Clause 33.06.07.02: Function Call Once},pdfsubject={Cornell notes},pdfauthor={}}

\begin{document}
\maketitle
\begin{CornellOverviewBox}{Overview}
\textbf{Classification:} Standard library - Normative\\[2pt]
\textbf{Learning objective:} Understand the rules, terminology, and semantic consequences associated with Function call\_once.
\end{CornellOverviewBox}\begin{CornellSummaryBox}{Summary}
{\sffamily\bfseries\color{Accent} Summary}\par\vspace{3pt}Section 33.6.7.2 focuses on Function call\_once. Apply the specified interface together with its mandates, effects, result conditions, exception behavior, and complexity constraints. When applying it, preserve the distinction between mandatory requirements, recommendations, permissions, and behavior deliberately left undefined, unspecified, or implementation-defined.
\end{CornellSummaryBox}

\end{document}
